% !TeX root = ../main.tex
% ---------------------------------------------------------------------------
%  Three-Dimensional Ground Modelling
%  New chapter. The ground model is the object every later chapter consumes,
%  and it is now built, updated and interrogated digitally. Implicit modelling,
%  geostatistics, machine learning and uncertainty representation are treated
%  here rather than scattered through the investigation and digital chapters.
%  Case material should be drawn from published EIS/EES documents, geotechnical
%  baseline reports and peer-reviewed papers -- not from news reporting.
% ---------------------------------------------------------------------------
\chapter{Three-Dimensional Ground Modelling}
\label{ch:ground-modelling}

Every subsequent chapter consumes a ground model. Empirical classification
needs domains; analytical solutions need a stratigraphy and a stress state;
numerical analysis needs geometry, material assignment and boundaries;
construction control needs a chainage-by-chainage expectation to compare
against what the face actually reveals. The model is therefore not a
by-product of investigation but the principal engineering deliverable of it.

What has changed is how the model is built. Hand-drawn sections interpolated
between boreholes have given way to implicit surfaces fitted to all available
data, updated as each new borehole, geophysical line or mapped face arrives,
and increasingly conditioned by statistical and machine-learning methods. The
gain is not draughting speed. It is that the model can be rebuilt cheaply
enough to be kept current, and that its uncertainty can be quantified rather
than implied by the confidence of the draughtsman.

\section{From data to ground model}
\label{sec:data-to-ground-model}

\subsection{Model types and their purposes}
\label{sec:model-types}
\scopenote{The geological model, the hydrogeological model, the geotechnical
ground model and the analysis model: distinct objects with distinct
tolerances, frequently conflated. Which decisions each is fit to support.}

\subsection{Data sources and their support volumes}
\label{sec:data-support}
\scopenote{Boreholes, mapped exposures, geophysical sections, face mapping and
machine data have very different spatial support. Combining them without
accounting for that difference produces a model more precise than the data
justifies.}

\subsection{Model resolution and decision scale}
\label{sec:model-resolution}
\scopenote{Resolution should follow the decision. Route selection needs
formation-scale geometry; face support classification needs metre-scale
structure. Modelling everything at the finest scale is expensive and
misleading.}

\section{Building the geometry}
\label{sec:building-geometry}

\subsection{Explicit and implicit modelling}
\label{sec:implicit-modelling}
\scopenote{Explicit construction from digitised sections against implicit
methods that fit a scalar field to contacts and orientation data. Why implicit
models can be regenerated on each data update, and where they produce
geologically implausible surfaces that require explicit control.}

\subsection{Structural and stratigraphic control}
\label{sec:structural-control-modelling}
\scopenote{Imposing geological rules --- stratigraphic order, fault
displacement, fold geometry, intrusive relationships --- so the interpolator
produces a geologically admissible result rather than a smooth one.}

\subsection{Faults, domains and discontinuity networks}
\label{sec:dfn}
\scopenote{Representing faults as modelled surfaces; discrete fracture network
generation from the orientation, spacing and persistence statistics of
Chapter~\ref{ch:structural-analysis}; conditioning networks to observed
intersections.}

\section{Populating the model}
\label{sec:populating-model}

\subsection{Geostatistical estimation and simulation}
\label{sec:geostatistics}
\scopenote{Variograms, kriging and conditional simulation applied to
geotechnical properties. The distinction that matters: estimation gives a
smoothed best guess, simulation gives equally probable realisations that
preserve variability --- and design against variability needs the second.}

\subsection{Machine learning for property and lithology prediction}
\label{sec:ml-ground-modelling}
\scopenote{Supervised classification of lithology and weathering grade from
downhole logs and geophysics; regression for strength and modulus; automated
core imagery interpretation. Training data volume, class imbalance, and
validation against withheld boreholes rather than against the training set.}

\subsection{Assigning parameters and domains}
\label{sec:domain-assignment}
\scopenote{Converting the populated model into geotechnical domains with
design parameters attached, and carrying the parameter derivation chain of
Chapter~\ref{ch:rock-mass-parameters} into the model rather than into a
separate spreadsheet.}

\section{Uncertainty}
\label{sec:model-uncertainty}

\subsection{Sources of model uncertainty}
\label{sec:model-uncertainty-sources}
\scopenote{Data density and distribution, interpretation ambiguity,
interpolation method and geological complexity. A model built from three
boreholes and one built from three hundred render identically.}

\subsection{Representing and communicating uncertainty}
\label{sec:uncertainty-representation}
\scopenote{Multiple realisations, probability-of-lithology volumes, distance
to nearest datum, and confidence classification along the alignment. Practical
means of stopping a rendered surface from being read as a survey.}

\subsection{Value of information}
\label{sec:value-of-information}
\scopenote{Using the model to decide where the next borehole should go:
targeting investigation at the uncertainty that actually changes a design
decision, rather than at uniform spacing.}

\section{Keeping the model current}
\label{sec:model-updating}

\subsection{Updating during investigation and construction}
\label{sec:model-update-cycle}
\scopenote{Bayesian and sequential updating as boreholes, probe holes,
geophysics and mapped faces arrive; face mapping and machine data as a
continuous data stream ahead of and behind the excavation.}

\subsection{Reconciling prediction with observation}
\label{sec:prediction-vs-observation}
\scopenote{Systematic comparison of predicted against encountered conditions,
chainage by chainage. This is the evidence base for baseline comparison
(Chapter~\ref{ch:delivery}) and for revising the classification ahead
of the face.}

\subsection{The model as a project record}
\label{sec:model-as-record}
\scopenote{Version control, provenance of every surface and parameter, and
handover of a model an operator can still use decades later
(Chapter~\ref{ch:monitoring-digital}).}
